\section{Polynomials $t_{j,p}(\boldsymbol{q})$}

\[
t_{j,p}(\boldsymbol{q}) = \sum_{d} f_{p,d}^{j}\boldsymbol{q}^{d}
\]

\begin{align*}
\epsilon_1 &= \nu_{1,0} + \nu_{1,4} + \nu_{5,1} \\
\epsilon_2 &= \nu_{2,0} \\
\epsilon_3 &= \nu_{3,0} + \nu_{5,1} + \nu_{6,2} \\
\epsilon_4 &= \nu_{4,0} \\
\epsilon_5 &= \nu_{1,1} + \nu_{1,3} + \nu_{3,1} + \nu_{5,0} + \nu_{5,2} + \nu_{6,1} \\
\epsilon_6 &= \nu_{3,2} + \nu_{5,1} + \nu_{6,0} \\
\epsilon_7 &= \nu_{7,0} \\
\epsilon_8 &= \nu_{8,0}
\end{align*}

\[
\begin{array}{c|c|c|c|c|c|c|c|c|c}
j & \text{Simple} & t_{j,1}(\boldsymbol{q}) & t_{j,2}(\boldsymbol{q}) & t_{j,3}(\boldsymbol{q}) & t_{j,4}(\boldsymbol{q}) & t_{j,5}(\boldsymbol{q}) & t_{j,6}(\boldsymbol{q}) & t_{j,7}(\boldsymbol{q}) & t_{j,8}(\boldsymbol{q}) \\\hline
1 & M(e,+) & 1 + \boldsymbol{q}^{4} & 0 & 0 & 0 & \boldsymbol{q} & 0 & 0 & 0 \\
2 & M(e,-) & 0 & 1 & 0 & 0 & 0 & 0 & 0 & 0 \\
3 & M(e,\rho) & 0 & 0 & 1 & 0 & \boldsymbol{q} & \boldsymbol{q}^{2} & 0 & 0 \\
4 & M(\sigma,+) & 0 & 0 & 0 & 1 & 0 & 0 & 0 & 0 \\
5 & M(\sigma,-) & \boldsymbol{q} + \boldsymbol{q}^{3} & 0 & \boldsymbol{q} & 0 & 1 + \boldsymbol{q}^{2} & \boldsymbol{q} & 0 & 0 \\
6 & M(\tau,0) & 0 & 0 & \boldsymbol{q}^{2} & 0 & \boldsymbol{q} & 1 & 0 & 0 \\
7 & M(\tau,1) & 0 & 0 & 0 & 0 & 0 & 0 & 1 & 0 \\
8 & M(\tau,2) & 0 & 0 & 0 & 0 & 0 & 0 & 0 & 1 \\
\end{array}
\]
